\documentclass[quick,pro]{debate}
\motion{AI的迅猛发展提升 / 降低了人类创作者存在的意义}
\stance{正方}{提升了}
\begin{document}
\debatetitle
\begin{thesisbox}机器能做的越多，“这是一个人做的”就越被区分、被看重，作品的分量也越落在人的立意上。\end{thesisbox}

\section{定义与判准（标准表述）}
\begin{fields}
\field{AI 的迅猛发展}{2022 年以来，机器生成文字、图像、音乐和声音的能力突然变强，而且人人可用。}
\field{人类创作者}{把创造文艺作品当成自己一部分的人，靠它吃饭的算，业余在写在画的也算。}
\field{存在的意义}{作用与重要性（汉典）。两层：这个世界还有多需要、多珍视由人来创作；创作还能给创作者自己多少东西。}
\field{判准}{和 AI 爆发之前比，由人来创作这件事整体变重还是变轻；需要和珍视两半都算，看多数创作者。\sub{我方要证明}{珍视这一半在变重，而且不只落在顶尖少数；立意在作品里的分量更重。}\sub{对方要证明}{需要和珍视加在一起，对多数创作者变轻了。}}
\field{承认线}{一部分执行型的订单确实在减少。但题目问的是意义，不是订单数。}
\field{偏向定义 30 秒口径}{意义是作用与重要性，被需要和被珍视两半都算。照“意义就是收入”，一个没有稿费的诗人没有存在的意义，这不是常识。}
\end{fields}

\section{我方论点}
\begin{tabularx}{\linewidth}{@{}L{5em}Y Y L{6.4em}@{}}\toprule
\thead{标签}&\thead{一句话}&\thead{核心数据 / 学理 / 案例}&\thead{出处}\\\midrule
\textbf{人味成了稀缺品}&机器内容无限以后，社会开始专门区分人和机器，并在能选的时候选人&数据：同一批 AI 画贴人类标签，四项评分全高；76\% 认为能分辨人或 AI 很重要。学理：原作价值来自一次创作行为。案例：法院只认人当作者；AI 内容须标识&Bellaiche 2023；Pew 2025；Newman \& Bloom 2012；Thaler 2025；网信办 2025\\
\textbf{立意回到人}&AI 接走画得像、写得顺，留给人的是想画什么、选哪一张&数据：5 万多用户 400 万件作品，产量 +25\%、获赞概率 +50\%，集中度下降；新手进步最大。案例：法院把 AI 图判给做选择的人；Randy Travis&Zhou \& Lee 2024；Doshi \& Hauser 2024；北京互联网法院 2023；CBS 2024\\
\bottomrule\end{tabularx}

\section{对方论点 → 一句话反驳}
\begin{tabularx}{\linewidth}{@{}Y Y Y@{}}\toprule
\thead{对方会说}&\thead{我方一句话回}&\thead{对方追问时}\\\midrule
不再被需要：订单收入在降，顶尖的也挡不住&订单少了，说明机器能做的活变多了，不说明人做的东西变轻了&“丢活不是事实？”→ 是，我方承认；丢一份活不等于人来创作变轻\\
手艺被掏空：人成了修图工&机器接走的是手上的重复，想画什么机器替不了&“视觉新颖度全降？”→ 平均在趋同，分数给了探索新想法的人\\
被区分不等于被需要&被认出来是被选择的前提；我方说珍视，不说价格&“标签带来订单吗？”→ 题目问意义，不问接单量\\
九成七的人盲听分不出&同一调查里八成要求标注，超过一半为此不安&“那是怕被骗”→ 怕被骗，就是在乎背后是不是人\\
保护说明濒危&意义是重要性，不是安全感；人只为看重的东西立法&“标识为了防虚假信息”→ 结果是“是不是人做的”成了公开事实\\
\bottomrule\end{tabularx}

\section{必问与必答}
\begin{points}{必问（对方不答就复述一次，不换题）}
\item 如果人的来源不重要，为什么法院、行业和七成以上的普通人，都坚持要把人和机器分开？（全场追问）
\item 贵方说的意义，包不包括被珍视？珍视那一半，贵方算过没有？
\item 北京那个 AI 图的案子，著作权判给了人还是 AI？
\end{points}
\begin{points}{必答（15 秒正面回答）}
\item 插画师丢了活，贵方承认吗？ → 承认执行型订单在减少；题目问意义，不问订单数。
\item 标签能带来订单吗？ → 被认出来是被选择的前提；我方说珍视，不说价格。
\item Randy Travis 是明星个例吧？ → 是个例，只说明机制；普遍性靠 5 万多名用户的数据和法院规则。
\end{points}

\section{对辩战场概要（三场）}
\begin{tabularx}{\linewidth}{@{}L{5.4em}Y Y Y@{}}\toprule
\thead{对辩}&\thead{开场问题}&\thead{目标}&\thead{收束语}\\\midrule
正三（开第一战场）&珍视那一半，贵方算过没有？&把战场定成“两半都算”&对方只量了按件计价的那一半\\
正一（中段收拢）&标识和法律挑不挑名气？&回应“珍视只给少数”，拉回看多数&规则对每个人一样，这就是多数\\
正二（结辩前）&被压价的是执行，还是想画什么？&守住“立意回到人”&机器接走的是手，接不走想说话的人\\
\bottomrule\end{tabularx}

\section{如果……就……}
\begin{contingency}
\ifthen{对方把意义定义成收入}{背口径：意义是作用与重要性，两半都算；没有稿费的诗人也有存在的意义。}
\ifthen{对方拿出没预料到的数据}{先问出处和样本，再问它量的是需要还是珍视、是多数还是少数。}
\ifthen{“人味稀缺”被打穿}{收缩到“立意回到人”，用法院规则说明作品归属看人的选择。}
\end{contingency}

各环节：质询全是 90 秒双边，不频繁打断；正二申论先回应反二；正一质询反三为结辩取材；正一结辩是最后一发，留约 40 字临场位。

\begin{recordcard}
\record{反一立论后}{反方立论的主干}{\opt{反方立论建立在「订单和收入的下降」上}\opt{反方立论建立在「作品里属于人的部分变少」上}\opt{反方立论建立在「观众分不出人和机器」上}\opt{反方立论的主干不在以上几条时}}{正二申论·拆反方立论}
\record{反二申论后}{反二申论主打什么}{\opt{反二申论主打「被区分不等于被需要」}\opt{反二申论主打「立意回到的是甲方」}\opt{反二申论主打「订单和收入在下降」}\opt{反二申论的重点不在以上几条时}}{正二申论·回应反二}
\record{一辩对辩后}{反方到此最有力的一点}{\opt{反方全场最有力的是「修图工」}\opt{反方全场最有力的是「订单和收入数据」}\opt{反方全场最有力的是「分不出人和机器」}\opt{反方最有力的一点不在以上几条时}}{正三申论·处理反方最强点}
\record{二辩对辩后}{全场主战场落在}{\opt{全场主战场落在「需要和珍视哪个算数」}\opt{全场主战场落在「修图工还是立意者」}\opt{主战场不清楚时}}{正一结辩·归纳分歧}
\record{二辩对辩后}{反方全场最有力的一点}{\opt{反方全场最有力的是「订单和收入数据」}\opt{反方全场最有力的是「修图工」}\opt{反方最有力的一点不在以上几条时}}{正一结辩·处理最强点}
\record{反一结辩后}{反一结辩收在哪里}{\opt{反一结辩收在「珍视只落在少数人身上」}\opt{反一结辩收在「用 AI 你变少，不用 AI 需要你的人变少」}\opt{反一结辩收在一个被替代的创作者身上}\opt{反一结辩的收束不在以上几条时}}{正一结辩·回应反一，再用临场位}
\record{全场}{对方回避了哪些问题}{（写下来）}{申论与结辩里点名}
\end{recordcard}
\end{document}
